% ============================================================
% SECTION 5 - THEORETICAL RESULTS
% ============================================================
\section{Theoretical Results}
\label{sec:theory}

This section contains the main theoretical contributions of the
paper.  The within-stratum assumptions for \texttt{EnbPI-S} are
stated in Section~\ref{sec:theory:assumptions}; the stratified
coverage bound (Theorem~\ref{thm:strat_coverage}) is presented in
Section~\ref{sec:theory:main} and the corresponding marginal result
in Section~\ref{sec:theory:marginal}; and
Section~\ref{sec:theory:dominance} derives the condition under which
\texttt{EnbPI-S} outperforms pooled EnbPI.  Proofs are in
Appendix~\ref{app:proofs}, except for the short ones given in place
here and in Section~\ref{sec:model:seasonal}.

%------------------------------------------------------------
\subsection{Within-Stratum Assumptions}
\label{sec:theory:assumptions}
%------------------------------------------------------------

The following two assumptions are the season-specific analogues
of Assumptions~\ref{ass:xu_a1} and~\ref{ass:xu_a2}.

\begin{assumption}[Within-stratum i.i.d.\ errors]
\label{ass:strat_a1}
The errors
$\{\epsilon_t : t \in \calT_{\sstar}\} \cup \{\epsilon_{T+1}\}$
are i.i.d.\ according to $F_{\sstar}$, which is Lipschitz
continuous with constant $L_{\sstar} > 0$.
\end{assumption}

\begin{assumption}[Within-stratum estimation quality]
\label{ass:strat_a2}
There exists a real sequence $\{\deltaTs\}_{T \geq 1}$
(implicitly indexed through $T_{\sstar} = T/\Sper$) such that
\begin{equation}
  \label{eq:strat_a2}
  \frac{1}{\Ts}
  \sum_{t \in \calT_{\sstar}}
  \bigl(\hat{f}_{-t}(x_t) - f(x_t)\bigr)^2
  \;\leq\; \deltaTs^2
  \qquad\text{and}\qquad
  \bigl|\hat{f}_{-(T+1)}(x_{T+1}) - f(x_{T+1})\bigr|
  \;\leq\; \deltaTs.
\end{equation}
\end{assumption}

\begin{remark}[Relationship between assumptions]
\label{rem:assumption_nesting}
Assumption~\ref{ass:strat_a1} follows directly from
Assumption~\ref{ass:seasonal_general}, which already states that the
errors of a given season are independent and identically distributed
according to a Lipschitz $F_{\sstar}$; the scale-mixture structure of
Assumption~\ref{ass:seasonal} is not needed, and under it
$F_{\sstar} = G(\cdot/\sigma_{\sstar})$.
Assumption~\ref{ass:strat_a2} is implied by
Assumption~\ref{ass:xu_a2}: if the average MSE over all $T$
training points is at most $\delta_T^2$, then the average
over any subset (including $\calT_{\sstar}$) is at most
$\Sper \cdot \delta_T^2$, so \eqref{eq:strat_a2} holds with
$\deltaTs \leq \sqrt{\Sper}\,\deltaT$.
\end{remark}

%------------------------------------------------------------
\subsection{Stratified Coverage Bound}
\label{sec:theory:main}
%------------------------------------------------------------

\begin{theorem}[\texttt{EnbPI-S} conditional coverage bound]
\label{thm:strat_coverage}
Let Assumptions~\ref{ass:strat_a1} and~\ref{ass:strat_a2} hold.
Then, for any $T \in \N$, $\alpha \in (0,1)$, and
$\beta \in [0,\alpha]$:
\begin{align}
  \label{eq:strat_bound}
  &\bigl|
    \Prob\!\bigl(
      Y_{T+1} \in \hat{C}_{T+1}^{\,\alpha,\textup{strat}}
      \mid X_{T+1} = x_{T+1}
    \bigr) - (1-\alpha)
  \bigr|
  \notag\\
  &\quad\leq\;
  12\sqrt{\frac{\log(16\,\Ts)}{\Ts}}
  \;+\;
  4(L_{\sstar}+1)\!\left(\deltaTs^{2/3} + \deltaTs\right).
\end{align}
Furthermore, if $\deltaTs \to 0$ as $\Ts \to \infty$
(which occurs whenever $T \to \infty$ for fixed $\Sper$),
both terms on the right-hand side of~\eqref{eq:strat_bound}
converge to zero, establishing \emph{asymptotic conditional
validity} for every season $\sstar \in \{1,\ldots,\Sper\}$.
\end{theorem}

\begin{proof}
  See Appendix~\ref{app:proofs:thm1s}.  The proof applies the
  argument of \citet[Theorem~1]{xu2023conformal} to the
  within-stratum sub-sequence of length $\Ts$.  The key steps
  are: (i) observe that the stratified empirical $p$-value
  $\hat{p}_{T+1}^{(\sstar)} = \hat{F}_{\Ts}^{(\sstar)}(\hat{\epsilon}_{T+1})$
  is defined identically to the pooled one but using only
  season-$\sstar$ residuals; (ii) the Probability Integral
  Transform gives $F_{\sstar}(\epsilon_{T+1}) \sim \mathrm{Unif}[0,1]$
  since $\epsilon_{T+1} \sim F_{\sstar}$
  (Assumption~\ref{ass:strat_a1}); (iii) the DKW inequality
  \citep{massart1990tight} applied to the $\Ts$ i.i.d.\
  within-stratum errors yields the first term; (iv) the LOO
  coupling \eqref{eq:coupling} and Assumption~\ref{ass:strat_a2}
  yield the second term.  No seasonal bias term appears because
  the stratified empirical CDF estimates $F_{\sstar}$ directly.
\end{proof}

\begin{remark}[Comparison with Theorem~\ref{thm:xu_original}]
\label{rem:compare_theorems}
Theorem~\ref{thm:strat_coverage} and Theorem~\ref{thm:xu_original},
the original result of \citet{xu2023conformal}, have the same
functional form, since both bound the conditional coverage gap by
the sum of a DKW-type sampling term and an estimation quality term.
They differ on two counts.  First, Theorem~\ref{thm:xu_original}
uses $T$, the total number of training observations, and $\deltaT$,
whereas Theorem~\ref{thm:strat_coverage} uses $\Ts = T/\Sper$ and
$\deltaTs$, so that the stratified bound converges at the slower
rate $O(\sqrt{\Sper}\cdot T^{-1/2})$, a consequence of the smaller
effective sample size.  Second, Theorem~\ref{thm:xu_original} is
\emph{only valid} when Assumption~\ref{ass:xu_a1} holds, that is,
when all errors share a common CDF, which fails under seasonal
heteroskedasticity; Proposition~\ref{prop:pooled_gap} adds the
seasonal bias $2\bsstar$ to the right-hand side of
Theorem~\ref{thm:xu_original}, and adjusts the sampling constant,
yielding a bound that is valid but does not converge to zero.
Theorem~\ref{thm:strat_coverage} has no such persistent term.
\end{remark}

\begin{remark}[Extension to dependent within-stratum errors]
\label{rem:mixing}
Assumption~\ref{ass:strat_a1} can be relaxed exactly as in
\citet[Corollaries~1--2]{xu2023conformal}: if within-stratum errors
follow a linear process, the DKW term becomes
$O(\log \Ts / \sqrt{\Ts})$; if they are strongly mixing with
summable mixing coefficients, it becomes
$O((\log \Ts)^{2/3} / \Ts^{1/3})$.
The within-stratum errors $\{\epsilon_t : t \in \calT_{\sstar}\}$
are spaced $\Sper$ time steps apart, so that, if the original
process is $\alpha_n$-mixing, the sub-sampled process has mixing
coefficients bounded by $\alpha_{\Sper}, \alpha_{2\Sper}, \ldots$
(Proposition~\ref{prop:submixing} in Appendix~\ref{app:proofs}).
For most economic time series $\Sper = 12$ is large enough that the
within-stratum errors are approximately i.i.d., which makes
Assumption~\ref{ass:strat_a1} a reasonable working hypothesis; the
autocorrelation function of the within-stratum LOO residuals can be
used as a diagnostic (cf.\ \citealt{xu2023conformal}).
\end{remark}

%------------------------------------------------------------
\subsection{Marginal Coverage}
\label{sec:theory:marginal}
%------------------------------------------------------------

\begin{corollary}[Marginal coverage within a season]
\label{cor:marginal}
Under the conditions of Theorem~\ref{thm:strat_coverage},
the same bound holds for the marginal coverage conditional
on the season:
\begin{equation}
  \label{eq:marginal_strat}
  \bigl|
    \Prob\!\bigl(
      Y_{T+1} \in \hat{C}_{T+1}^{\,\alpha,\textup{strat}}
      \mid s(T+1) = \sstar
    \bigr) - (1-\alpha)
  \bigr|
  \;\leq\;
  12\sqrt{\frac{\log(16\,\Ts)}{\Ts}}
  \;+\;
  4(L_{\sstar}+1)\!\left(\deltaTs^{2/3} + \deltaTs\right).
\end{equation}
\end{corollary}

\begin{proof}
Apply the law of total expectation to
Theorem~\ref{thm:strat_coverage}, integrating over the conditional
distribution of $X_{T+1}$ given $s(T+1) = \sstar$.  The bound
is the same since the right-hand side of~\eqref{eq:strat_bound}
does not depend on $x_{T+1}$.
\end{proof}

%------------------------------------------------------------
\subsection{When Does Stratification Dominate Pooling?}
\label{sec:theory:dominance}
%------------------------------------------------------------

Comparing the asymptotic behaviour of the two bounds, the
structural result is immediate.

\begin{proposition}[Asymptotic dominance of stratification]
\label{prop:strat_dominance}
Suppose Assumptions~\ref{ass:strat_a1}--\ref{ass:strat_a2}
and~\ref{ass:seasonal} hold, and that $\deltaTs \to 0$ and
$\deltaT \to 0$ as $T \to \infty$.  Then:
\begin{enumerate}
  \item[\textup{(i)}]
    The right-hand side of the pooled bound
    (Proposition~\ref{prop:pooled_gap}) does not converge to zero:
    \[
      \lim_{T \to \infty}
      \Bigl[
        c_0\sqrt{\tfrac{\log(16T)}{T}}
        + 4(L+1)\!\left(\deltaT^{2/3}+\deltaT\right)
        + 2\bsstar
      \Bigr]
      \;=\; 2\bsstar.
    \]
    Hence the pooled coverage bound is bounded below by
    $2\bsstar > 0$ for all $T$: no amount of
    additional data drives the bound to zero.

  \item[\textup{(ii)}]
    The right-hand side of the stratified bound
    (Theorem~\ref{thm:strat_coverage}) converges to zero, and
    therefore:
    \[
      \lim_{T \to \infty}
      \bigl|
        \Prob(Y_{T+1} \in \hat{C}_{T+1}^{\alpha,\textup{strat}}
              \mid X_{T+1} = x_{T+1})
        - (1-\alpha)
      \bigr|
      \;=\; 0.
    \]

  \item[\textup{(iii)}]
    Consequently, whenever $\bsstar > 0$, the pooled bound
    remains bounded below by
    $2\bsstar > 0$ for all $T$, while the stratified guarantee
    converges to zero.  By Lemma~\ref{lem:bias_properties}(iii),
    $\bsstar > 0$ whenever $G$ is strictly increasing and season
    $\sstar$ has the strictly largest or smallest seasonal scale
    (in particular, for a genuinely high-volatility season).
    Stratification therefore provides a strictly tighter theoretical
    guarantee than pooling whenever such seasonal heteroskedasticity
    is present.
\end{enumerate}
\end{proposition}

\begin{proof}
  Part~(i): the sampling term
  $c_0\sqrt{\log(16T)/T} \to 0$ and the estimation term
  $4(L+1)(\deltaT^{2/3}+\deltaT) \to 0$ by hypothesis; the bias
  term $2\bsstar$ is independent of $T$
  (Lemma~\ref{lem:bias_properties}(ii)), so the limit of the
  right-hand side of~\eqref{eq:pooled_gap} is $2\bsstar$.
  Part~(ii) follows from Theorem~\ref{thm:strat_coverage}: both
  terms on the right of~\eqref{eq:strat_bound} vanish as
  $\Ts \to \infty$ (which occurs as $T\to\infty$ for fixed $\Sper$).
  Part~(iii) combines (i) and (ii) with
  Lemma~\ref{lem:bias_properties}(iii).
\end{proof}

Proposition~\ref{prop:strat_dominance} is asymptotic, and for finite
$T$ there is a genuine trade-off, since stratification removes the
persistent bias $2\bsstar$ but increases the sampling term by a
factor of $\sqrt{\Sper}$, from $O(\sqrt{\log T / T})$ to
$O(\sqrt{S\log T / T})$.  Stratification is preferred over pooling
for a given season $\sstar$ when the bias dominates the additional
sampling variance.  Since the pooled sampling and estimation terms
are nonnegative, and $L \geq L_{\sstar}$, a conservative sufficient
condition for the stratified bound~\eqref{eq:strat_bound} to be
smaller than the pooled bound~\eqref{eq:pooled_gap} is:
\begin{equation}
  \label{eq:dominance_condition}
  2\bsstar \;>\;
  12\sqrt{\frac{\log(16\Ts)}{\Ts}}
  + 4(L_{\sstar}+1)
  \left[\bigl(\deltaTs^{2/3} + \deltaTs\bigr)
        - \bigl(\deltaT^{2/3} + \deltaT\bigr)\right],
\end{equation}
where the right-hand side is positive whenever
$\deltaTs \geq \deltaT$, which is the typical case, given that the
stratum uses fewer observations.

\begin{remark}[Practical applicability of condition~\eqref{eq:dominance_condition}]
\label{rem:dominance_practical}
For $\Sper = 12$ and $T = 120$ (10 years of monthly data),
$\Ts = 10$ and the first term on the right-hand side of
\eqref{eq:dominance_condition} is
\[
  12\sqrt{\frac{\log(16\Ts)}{\Ts}}
  = 12\sqrt{\frac{\ln(160)}{10}}
  \approx 12 \times 0.713 \approx 8.55,
\]
which alone exceeds $2$, twice the maximum possible value of
$\bsstar$, so that the condition cannot be satisfied at $T = 120$.
The difficulty is not confined to that sample size.  Ignoring the
estimation terms, \eqref{eq:dominance_condition} requires
$\Ts > 36\,\bsstar^{-2}\log(16\Ts)$, so even the largest conceivable
bias, $\bsstar = 1$, would call for $\Ts \approx 300$ observations
per season, and the value $b_H \approx 0.12$ of the Gaussian example
in Section~\ref{sec:model:seasonal} would call for $\Ts$ of the order
of $10^{4}$.  No economic time series comes close.  Condition
\eqref{eq:dominance_condition} should therefore be read as a
qualitative statement about which term dominates, and not as a
threshold to be checked in an application; the same reservation
applies to the bound of \citet[Theorem~1]{xu2023conformal} itself,
whose right-hand side exceeds one at every sample size used in this
paper.  What the finite-sample comparison can be based on is
simulation, and Section~\ref{sec:simulations} shows that
\texttt{EnbPI-S} achieves better high-volatility conditional coverage
than pooled EnbPI from $T = 240$ months onwards
(Table~\ref{tab:e1_coverage_T240}), in the direction predicted by
Proposition~\ref{prop:strat_dominance}.
\end{remark}

%------------------------------------------------------------
\subsection{Width Bound for \texttt{EnbPI-S}}
\label{sec:theory:width}
%------------------------------------------------------------

The following result bounds the gap in interval width between
\texttt{EnbPI-S} and the oracle interval $C_{T+1}^{\alpha}$.

\begin{theorem}[\texttt{EnbPI-S} width gap bound]
\label{thm:width_bound}
Under Assumptions~\ref{ass:strat_a1}--\ref{ass:strat_a2},
further assume $F_{\sstar}^{-1}$ is Lipschitz with constant
$K_{\sstar} > 0$.  With probability at least
$1 - \sqrt{\log(16\Ts)/\Ts}$:
\begin{equation}
  \label{eq:width_bound}
  \mathcal{D}_{\sstar}(T)
  \;\leq\;
  \deltaTs
  + \frac{\alpha K'_{\sstar}}{m}
  + 2(K_{\sstar} + M_{\sstar})
  \left(
    3\sqrt{\frac{\log(16\Ts)}{\Ts}}
    + (L_{\sstar}+1)\!\left(\deltaTs^{2/3}+\deltaTs\right)
  \right),
\end{equation}
where:
\begin{itemize}
  \item $\mathcal{D}_{\sstar}(T) = \lambda\bigl(\hat{C}_{T+1}^{\alpha,\textup{strat}}
    \triangle C_{T+1}^{\alpha}\bigr)$ is the Lebesgue measure
    $\lambda$ of the symmetric difference between the stratified
    and oracle intervals;
  \item $m$ is the number of line-search grids for
    $\hat{\beta}^{(\sstar)}$ \eqref{eq:betahat_strat};
  \item $K'_{\sstar} = \max_{j=1,\ldots,\Ts-1}
    (\hat{\epsilon}_{(j+1)}^{(\sstar)} - \hat{\epsilon}_{(j)}^{(\sstar)})$
    is the largest gap between consecutive \emph{sorted}
    within-stratum LOO residuals
    ($\hat{\epsilon}_{(1)}^{(\sstar)} \leq \cdots \leq
      \hat{\epsilon}_{(\Ts)}^{(\sstar)}$);
  \item $M_{\sstar} = 2K_{\sstar}/|1/K'_{\sstar} - L_{\sstar}|$
    (well-defined provided $K'_{\sstar} \neq 1/L_{\sstar}$; since
    $K'_{\sstar}$ is a data-dependent quantity, this condition holds
    almost surely).
\end{itemize}
\end{theorem}

\begin{proof}
  See Appendix~\ref{app:proofs:thm2s}.  The proof mirrors
  \citet[Theorem~3]{xu2023conformal}, replacing all global
  quantities ($T$, $\deltaT$, etc.) with their stratum-$\sstar$
  counterparts ($\Ts$, $\deltaTs$, etc.).
\end{proof}

\begin{remark}
  Theorem~\ref{thm:width_bound} implies that, as $T \to \infty$
  and $\deltaTs \to 0$, the stratified interval
  $\hat{C}_{T+1}^{\alpha,\textup{strat}}$ converges in width
  to the oracle $C_{T+1}^{\alpha}$ for season $\sstar$.  The
  convergence rate is determined by $\Ts = T/\Sper$, with the
  same functional form as in the non-seasonal case.
\end{remark}

%------------------------------------------------------------
\subsection{The Finite-Sample Guarantee of the Conformal Construction}
\label{sec:theory:finite}
%------------------------------------------------------------

Theorem~\ref{thm:strat_coverage} is asymptotic, and its right-hand side
exceeds one at every sample size this paper works with.  The
construction that Section~\ref{sec:algorithm:defects} recommends admits
a stronger and much simpler statement, which is what the experiments of
Section~\ref{sec:simulations:e6constr} actually measure, and it is
worth recording because the argument is two lines and the result is
exact.

\begin{proposition}[Exact validity of the conformal endpoints]
\label{prop:finite_sample}
Let $\epsilon_1, \ldots, \epsilon_{n}, \epsilon_{n+1}$ be exchangeable
with a continuous joint distribution, write
$\epsilon_{(1)} \leq \cdots \leq \epsilon_{(n)}$ for the order
statistics of the first $n$, and set
$k_{\text{lo}} = \lfloor (n+1)\alpha/2 \rfloor$ and
$k_{\text{hi}} = \lceil (n+1)(1-\alpha/2) \rceil$.  If
$1 \leq k_{\text{lo}}$ and $k_{\text{hi}} \leq n$, then
\begin{equation}
  \label{eq:exact_coverage}
  1 - \alpha
  \;\leq\;
  \Prob\bigl(\epsilon_{(k_{\text{lo}})} \leq \epsilon_{n+1}
             \leq \epsilon_{(k_{\text{hi}})}\bigr)
  \;\leq\;
  1 - \alpha + \frac{2}{n+1}.
\end{equation}
Both index conditions hold if and only if $n \geq 2/\alpha - 1$, and
the two indices are interior, $k_{\text{lo}} \geq 2$ and
$k_{\text{hi}} \leq n-1$, if and only if $n \geq 4/\alpha - 1$.
\end{proposition}

\begin{proof}
By exchangeability and continuity the rank
$R = 1 + \#\{i \leq n : \epsilon_i < \epsilon_{n+1}\}$ is uniform on
$\{1,\ldots,n+1\}$.  Since
$\epsilon_{n+1} < \epsilon_{(k)} \iff R \leq k$, we have
$\Prob(\epsilon_{n+1} < \epsilon_{(k_{\text{lo}})})
= k_{\text{lo}}/(n+1) \leq \alpha/2$ because
$k_{\text{lo}} \leq (n+1)\alpha/2$, and
$\Prob(\epsilon_{n+1} > \epsilon_{(k_{\text{hi}})})
= (n+1-k_{\text{hi}})/(n+1) \leq \alpha/2$ because
$k_{\text{hi}} \geq (n+1)(1-\alpha/2)$.  Adding the two gives the lower
bound in~\eqref{eq:exact_coverage}; the upper bound follows because
each floor and each ceiling moves its index by less than one, so each
tail probability is at least $\alpha/2 - 1/(n+1)$.  For the index
conditions, $k_{\text{lo}} \geq 1$ requires
$(n+1)\alpha/2 \geq 1$ and $k_{\text{hi}} \leq n$ requires
$(n+1)(1-\alpha/2) \leq n$, and both reduce to
$n \geq 2/\alpha - 1$; the interior versions replace the right-hand
sides by $2$ and $n-1$ and reduce to $n \geq 4/\alpha - 1$.
\end{proof}

\begin{remark}[What the proposition does and does not deliver here]
\label{rem:finite_caveat}
Applied to a season buffer, Proposition~\ref{prop:finite_sample} says
that the interval of Section~\ref{sec:algorithm:defects} covers at
least $1-\alpha$ whatever $\Ts$ is, provided only that
$\Ts \geq 2/\alpha - 1$ and that the season's errors are exchangeable
with the one being predicted.  That is a stronger statement than
Theorem~\ref{thm:strat_coverage} and it does not weaken as the buffer
shortens; the upper bound $1-\alpha + 2/(\Ts+1)$ is what makes the
interval conservative rather than invalid at $\Ts = 20$, and it is
$0.995$ there, which is why the scheme overshoots in
Section~\ref{sec:simulations:e7schemes} and on the export application.

Two things stand between the proposition and the algorithm, and both
are inherited rather than introduced.  The buffer holds leave-one-out
residuals rather than errors, so exchangeability with the test residual
holds only up to the estimation error that
Assumption~\ref{ass:strat_a2} controls, and the coupling of
Lemma~\ref{lem:coupling} converts that into the same
$\deltaTs$ terms that already appear in
Theorem~\ref{thm:strat_coverage}.  And the sliding-window update makes
the buffer depend on the recent past, so exchangeability is exact only
within a season under Assumption~\ref{ass:strat_a1}.  We therefore
read~\eqref{eq:exact_coverage} as the explanation of what
Section~\ref{sec:simulations:e6constr} measures rather than as a
theorem about the deployed algorithm, and we note that the measured
coverage, $0.904$ at $n = 20$ and $0.901$ at $n = 500$, sits where the
proposition says it should.
\end{remark}

%------------------------------------------------------------
\subsection{Coverage Along the Family}
\label{sec:theory:family}
%------------------------------------------------------------

The family~\eqref{eq:family} interpolates between two estimators of the
same quantile, one of which is consistent for it and one of which is not
unless Assumption~\ref{ass:seasonal} holds.  The following bound says what
that costs.

\begin{proposition}[Coverage of the interpolated interval]
\label{prop:family_coverage}
Let Assumptions~\ref{ass:strat_a1} and~\ref{ass:strat_a2} hold and let
$F_{\sstar}$ be Lipschitz with constant $L_{\sstar}$.  Write
$\hat{C}_{T+1}^{\alpha,\lambda}$ for the interval built from
$\hat{\theta}^{\lambda}$ in~\eqref{eq:family} at the same quantile levels
as the stratified interval, and let
\[
  d_{\text{lo}}
  = \hat{\theta}^{N}(\beta) - \hat{\theta}^{S}(\beta),
  \qquad
  d_{\text{hi}}
  = \hat{\theta}^{N}(1-\alpha+\beta) - \hat{\theta}^{S}(1-\alpha+\beta).
\]
Then, for every $\lambda \in [0,1]$,
\begin{equation}
  \label{eq:family_bound}
  \bigl|
    \Prob\bigl(Y_{T+1} \in \hat{C}_{T+1}^{\alpha,\lambda}
               \mid X_{T+1} = x_{T+1}\bigr) - (1-\alpha)
  \bigr|
  \;\leq\;
  12\sqrt{\frac{\log(16\Ts)}{\Ts}}
  + 4(L_{\sstar}+1)\bigl(\deltaTs^{2/3} + \deltaTs\bigr)
  + (1-\lambda) L_{\sstar}\bigl(|d_{\text{lo}}| + |d_{\text{hi}}|\bigr).
\end{equation}
Moreover $|d_{\text{lo}}| + |d_{\text{hi}}| \leq
2\max_p|\Delta_{\sstar}(p)| + O_p(\Ts^{-1/2})$, so that:
\textup{(i)} at $\lambda = 1$ the bound is exactly that of
Theorem~\ref{thm:strat_coverage}; \textup{(ii)} if
Assumption~\ref{ass:seasonal} holds, so that
$\Delta_{\sstar} \equiv 0$, the third term is $O_p(\Ts^{-1/2})$ for every
fixed $\lambda$ and the whole family is asymptotically valid;
\textup{(iii)} if $\Delta_{\sstar} \not\equiv 0$, asymptotic validity
requires $\lambda \to 1$.
\end{proposition}

\begin{proof}
Conditional on the calibration buffers, the two intervals have the same
centre and endpoints differing by $(1-\lambda)d_{\text{lo}}$ and
$(1-\lambda)d_{\text{hi}}$.  For any $a \leq b$ and $a' \leq b'$,
$|[F_{\sstar}(b) - F_{\sstar}(a)] - [F_{\sstar}(b') - F_{\sstar}(a')]|
\leq L_{\sstar}(|b-b'| + |a-a'|)$ by Lipschitz continuity, so the two
conditional coverage probabilities differ by at most
$(1-\lambda)L_{\sstar}(|d_{\text{lo}}| + |d_{\text{hi}}|)$.  Adding this
to the bound of Theorem~\ref{thm:strat_coverage} and integrating over the
buffers gives~\eqref{eq:family_bound}.  The statement about
$d$ follows from writing $d = \Delta_{\sstar} + (\hat{\theta}^{N} -
\theta^{N}) - (\hat{\theta}^{S} - \theta^{S})$ and applying the quantile
central limit theorem to each estimation error.
\end{proof}

\begin{remark}[Where the interior of the family is worth visiting]
\label{rem:optimal_lambda}
Bound~\eqref{eq:family_bound} is monotone in $\lambda$ and therefore
recommends $\lambda = 1$, which is an artefact of bounding a bias
without crediting the variance it buys.  The mean squared error of
$\hat{\theta}^{\lambda}(p)$ tells the other half.  Writing
$\Delta = \Delta_{\sstar}(p)$, $v_S = \Var(\hat{\theta}^{S})$,
$v_N = \Var(\hat{\theta}^{N})$ and $c$ for their covariance,
\[
  \mathrm{MSE}(\lambda)
  = (1-\lambda)^2\Delta^2 + \lambda^2 v_S + (1-\lambda)^2 v_N
    + 2\lambda(1-\lambda)c ,
  \qquad
  \lambda^{*} = \frac{\Delta^2 + v_N - c}
                     {\Delta^2 + v_S + v_N - 2c},
\]
the denominator being $\Delta^2 + \Var(\hat{\theta}^{S} -
\hat{\theta}^{N})$ and the second derivative twice that, so
$\lambda^{*}$ minimises.  Both clips can bind: $\lambda^{*} < 0$ whenever
$c > \Delta^2 + v_N$, which at $\Delta = 0$ happens as soon as the
correlation between the two estimators exceeds $\sqrt{v_N/v_S}$, and the
two are computed from overlapping data.  Under
Assumption~\ref{ass:seasonal}, $\Delta = 0$ and $\lambda^{*}$ settles at a
constant determined by the ratio $v_N/v_S$ and that correlation, neither
tending to zero nor to one.  When $\Delta \neq 0$ is fixed,
$1 - \lambda^{*} = (v_S - c)/(\Delta^2 + v_S + v_N - 2c) =
O(1/(\Ts\Delta^2))$, so the shrinkage bias is $O(1/(\Ts\Delta))$, of
smaller order than the $O(\Ts^{-1/2})$ sampling error, and the
optimally weighted estimator has the same first-order limit law as the
stratified one while strictly improving finite-sample mean squared error.
That equivalence is pointwise in $\Delta$ and informative only while
$\Ts\Delta^2 \to \infty$: along sequences with
$\Delta \asymp \Ts^{-1/2}$, which is exactly the case of a shape
discrepancy comparable to sampling noise, $1 - \lambda^{*}$ stays bounded
away from zero and the shrinkage bias is of the same order as the standard
error.
\end{remark}

\begin{remark}[What is not claimed for the data-driven asymmetry]
\label{rem:beta_honesty}
The same scruple applies one level up, and the paper is explicit about
it because its own experiments measure the consequence.
Theorem~\ref{thm:strat_coverage} is stated for a \emph{fixed}
$\beta \in [0,\alpha]$, and its proof uses that fixed value at the
point where the Probability Integral Transform is applied
(Appendix~\ref{app:proofs:thm1s}, Step~5).  The interval actually
deployed in~\eqref{eq:enbpis_interval} uses
$\hat{\beta}^{(\sstar)}$, the minimiser of the width over a grid
computed from the same buffer, so the theorem does not cover it.  This
is inherited from \citet{xu2023conformal} rather than introduced here,
but Section~\ref{sec:simulations:e6constr} shows that at the buffer
lengths of interest the gap is not a formality: choosing the asymmetry
from the buffer costs between four and six percentage points of
coverage.  A bound that covered $\hat{\beta}^{(\sstar)}$ would have to
pay for the selection, and we do not have one; what we do instead is
fix $\beta = \alpha/2$ when the buffer is short, which restores the
finite-sample guarantee the fixed-$\beta$ theory describes.
\end{remark}

\begin{remark}[What is not claimed for the data-driven weight]
\label{rem:lambda_honesty}
Remark~\ref{rem:optimal_lambda} concerns a deterministic $\lambda$.  The
weight actually used in Section~\ref{sec:simulations} is chosen from the
same residuals that form the buffer, so $\hat{\theta}^{\hat\lambda}$ is a
nonlinear function of the calibration sample and neither the mean squared
error calculation above nor the conformal exchangeability argument
transfers to it.  We therefore claim nothing finite-sample for it and
document its behaviour by simulation.  A moment estimator of $\lambda^{*}$
built from a bootstrap of $(v_S, v_N, c)$ was tried and abandoned:
Section~\ref{sec:simulations:e6constr} reports why.
\end{remark}
